\documentclass[10pt,a4paper]{amsart}
\usepackage[margin=19mm]{geometry}
\usepackage{amsmath,amssymb,mathtools}
\usepackage[scr=rsfs,cal=euler]{mathalfa}
\usepackage{xfrac}
\usepackage[unicode,hidelinks,bookmarks=false]{hyperref}
\newcommand{\ie}{i.e.\ }
\newcommand{\cf}{cf.\ }
\newcommand{\resp}{respectively\ }
\newcommand{\Subsection}{\subsection}
\providecommand{\nobreakdash}{\nobreak\hbox{-}}
\setlength{\emergencystretch}{2em}
\multlinegap0pt
\usepackage{cancel}
\newtheorem{proposition}{Proposition}
\title{Moments for generalizations of a coin flip game}
\author{Jia Huang}
\date{Source: arXiv:2605.19904v1 (CC BY 4.0)}
\begin{document}
\maketitle

\noindent\emph{Source and licence.} Moments for generalizations of a coin flip game. Version arXiv:2605.19904v1 (CC BY 4.0), distributed by arXiv under the Creative Commons Attribution 4.0 International licence, http://creativecommons.org/licenses/by/4.0/, as recorded in the arXiv licence metadata for this exact version and shown on the abstract page https://arxiv.org/abs/2605.19904v1. Full licence notice: \texttt{licenses/arxiv-2605.19904-CC-BY-4.0.txt}.\par\smallskip

\noindent\textit{Editorial scope.} Counted unit: the article's proposition prop:Carlitz (source0.tex lines 179--184). The article's complete original proof of it is delivered with it (source0.tex lines 186--203, one \texttt{proof} environment of 18 lines). No mathematics is written, repaired or re-derived: the statement and the proof are the article's own lines, in the article's own notation, and no retained line is substituted or re-flowed.  No further source-local context is needed: every cross-reference of every retained line resolves inside the two delivered spans.  Nothing was omitted from any retained span, so the package carries no omission record.  No retained line contains a citation, so the wrapper carries no bibliography and no bibliography entry.  No retained line contains a figure environment, an image-inclusion command or drawing code, so no image is delivered and the package declares no asset dependency.  Editorial formatting only, in addition: an AMS \texttt{amsart} preamble that restates the article's own declarations and macros used by the retained lines, namely the environments \texttt{proposition} on separate counters, together with the standard packages those lines need.  The retained environments print in this document as Proposition 1 in order of appearance, whereas the article prints the same items with the numbering of its own full document.  The complete native source of the article as distributed in the arXiv e-print for this exact version is bundled unchanged as \texttt{sources/200952/source0.tex}.
\begin{proposition}\label{prop:Carlitz}
For any integers $n,k\ge0$, we have
\[ 
\sum_{d=0}^k d^n x^d = \sum_{i=0}^n \frac{e_{n,i} x^i - e_{n,i}^k x^{k+i+1}}{(1-x)^{n+1}}.
\]
\end{proposition}

\begin{proof}
Fix an arbitrary integer $k\ge0$.
We prove the desired identity by induction on $n$.
It is trivial when $n=0$. 
Assume the result holds for an arbitrary $n\ge0$.
Applying the operator $x\frac{d}{dx}$ gives
\begin{align*}
\sum_{d=0}^k d^{n+1} x^d =& \sum_{i=0}^n \frac{ (1-x) 
\left( ie_{n,i} x^i - (k+i+1) e_{n,i}^k x^{k+i+1} \right) 
+ (n+1) \left( e_{n,i} x^{i+1} - e_{n,i}^k x^{k+i+2} \right) }
{(1-x)^{n+2}} \\
=& \sum_{i=0}^{n+1} \frac{ (ie_{n,i} + (n+2-i)e_{n,i-1})x^i 
- ( (k+i+1)e_{n,i}^k + (n+1-k-i) e_{n,i-1}^k ) x^{k+i+1} } 
{(1-x)^{n+2}} \\
=& \sum_{i=0}^{n+1} \frac{ e_{n+1,i} x^i 
-  e_{n+1,i}^k x^{k+i+1} } {(1-x)^{n+2}}. \qedhere
\end{align*}
\end{proof}

\end{document}
